\section{Methodology}
\label{sec:method}

The system pipeline is: (i) train six independent components; (ii) tune per-model
decision thresholds on the validation set; (iii) search for the best ensemble fusion
strategy. Figure~\ref{fig:pipeline} provides an overview.

\begin{figure}[htbp]
\centering
\small
\renewcommand{\arraystretch}{1.2}
\begin{tabular}{c}
\textbf{Raw JSONL} (train / validation / test) \\
$\downarrow$ \\
\textbf{Preprocessing \& Stratified Split} \\
$\downarrow$ \\
\begin{tabular}{|c|c|c|c|c|c|}
\hline
Greek & XLM-R & XLM-R & Hybrid & Dictionary & NER \\
BERT & Large & Base & IR & Baseline & +EL \\
\hline
\end{tabular}
\\
$\downarrow$ \\
\textbf{Per-model threshold tuning} (validation set) \\
$\downarrow$ \\
\textbf{Metaheuristic Ensemble} (11 strategies + 3 compositions) \\
$\downarrow$ \\
\textbf{Final Predictions}
\end{tabular}
\caption{High-level system pipeline.}
\label{fig:pipeline}
\end{figure}

\subsection{Greek BERT Multi-Label Classifier}
\label{subsec:bert}

\subsubsection{Architecture.}
We fine-tune \texttt{nlpaueb/bert-base-greek-uncased-v1} (Greek BERT)~\cite{koutsikakis2020greekbert}
with a two-layer MLP head: mean+CLS concatenation pooling (1536-dim) $\to$ Dropout(0.2)
$\to$ Linear($1536\!\to\!768$, GELU) $\to$ Dropout(0.2) $\to$ Linear($768\!\to\!115$).
Concatenating mean-pooled hidden states and the \texttt{[CLS]} vector outperformed plain
CLS extraction by $\sim$0.8 percentage points in our ablation, as diagnostic evidence
is distributed across the full document. Max sequence length: 256 tokens.

\subsubsection{Training.}
Table~\ref{tab:bert_params} lists the final hyperparameters. ASL~\cite{ridnik2021asl} replaced binary cross-entropy (BCE) after Phase~2, raising
validation micro-F1 from $\approx$0.74 to 0.788 by suppressing easy negatives.
ASL applies asymmetric focusing to positive and negative terms:
\begin{equation}
\mathcal{L}_{\text{ASL}} = \sum_y \begin{cases}
  (1-p)^{\gamma^+} \log p & y=1 \\
  (p_m)^{\gamma^-} \log(1-p_m) & y=0
\end{cases}
\end{equation}
where $p_m = \max(p - m, 0)$ shifts and clips easy negatives, $\gamma^+=1$, $\gamma^-=4$, $m=0.05$.
\textbf{LLRD} (factor 0.90) updates lower encoder layers at $\text{LR}\times0.90^l$,
preserving pre-trained Greek representations while adapting upper layers to the clinical
domain. The model evolved across four phases:
BCE baseline (0.74) $\to$ ASL (0.788) $\to$ LLRD + MLP (0.811 at $t{=}0.5$) $\to$ final P4 tuning
(0.7984 at $t{=}0.6$, 0.8062 after per-label threshold tuning).
The drop from 0.811 to 0.7984 reflects the stricter active threshold ($t{=}0.6$ vs $t{=}0.5$); per-label tuning then recovers and exceeds the prior peak.

\begin{table}[h]
\centering
\caption{Greek BERT final hyperparameters.}
\label{tab:bert_params}
\begin{tabular}{ll}
\toprule
Hyperparameter & Value \\
\midrule
Max length / batch & 256 tok / 32 (16 + grad.\ accum.\ 2) \\
Learning rate / LLRD & $3\times10^{-5}$ / 0.90 per layer \\
Warmup / weight decay & 6\% / 0.01 \\
Loss & ASL ($\gamma^-=4$, $\gamma^+=1$, clip $=0.05$) \\
Epochs / patience & 30 / 7 (BF16) \\
Pooling / head & mean+CLS concat / MLP 1536$\to$768$\to$115 \\
Dropout & 0.2 \\
Active eval threshold & 0.6 \\
\bottomrule
\end{tabular}
\end{table}

\subsubsection{Threshold tuning.}
A two-pass sweep (step 0.01) selects $t_l^{*} = \arg\max_{t \in [0.45,0.75]} F1_l(t;\mathcal{D}_{\text{val}})$ per label independently. Pass~2 accepts per-label values only for labels with $\geq$15 positive validation examples and $\Delta$F1 $>0.0015$. This adds $\approx$0.8 percentage points over the active threshold ($t{=}0.6$) and $\approx$3.7 percentage points over a fixed $t{=}0.5$ baseline, bringing the tuned validation micro-F1 to 0.8062.

\subsection{XLM-RoBERTa Classifiers}
\label{subsec:xlmr}

XLM-RoBERTa's 250K multilingual vocabulary handles Greek--Latin code-switching better
than Greek BERT's 35K vocabulary.

\subsubsection{Long-document strategy.}
A sliding window (stride 256) handles documents exceeding 512 tokens. At inference, chunk logits are fused as $\hat{y} = 0.5\,\bar{y}_{\text{mean}} + 0.5\,y_{\text{max}}$, balancing average evidence with the most confident chunk signal.

\subsubsection{Large variant.}
ZLPR loss, effective batch 32 (batch 2, grad.\ accum.\ $\times$16), LR $2\times10^{-5}$,
weight decay 0.05, dropout 0.15, 40 epochs (patience 5). Validation F1 ceiling
0.7538; slower convergence and memory constraints (560M parameters) led us to
focus development on Greek BERT.

\subsubsection{Base variant.}
Implements MedicalModelWithDescriptions: the \texttt{[CLS]} representation
passes through 5 parallel dropout layers ($p=0.1i$, $i\in\{1\ldots5\}$) whose logits
are averaged (multi-sample dropout). A semantic anchoring term adds
cosine similarity between the document embedding and frozen ICD-10 description
embeddings, scaled by a learnable scalar $\alpha$:
\begin{equation}
\hat{y} = W\mathbf{h} + \alpha\cdot\tfrac{\mathbf{h}}{\|\mathbf{h}\|}\mathbf{D}^\top
\end{equation}
where rows of $\mathbf{D}$ are unit-normalised (frozen MiniLM embeddings of Greek ICD-10 descriptions), so the term is a true cosine similarity.
$\alpha$ converges from 0.099 to $\approx$0.12, confirming a consistent but small
semantic correction. Training: AdamW LR $2\times10^{-5}$, early stopping (patience 4),
per-class positive weights (ceiling 20). Post-processing: dynamic threshold tuning +
ICD-10 hierarchy enforcement + co-occurrence rules. Validation micro-F1: \textbf{0.8101}.
Despite using a smaller backbone, the Base model outperforms the Large variant
(0.8101 vs.\ 0.7538). The primary factors are the semantic anchoring mechanism and
task-specific post-processing absent in Large; GPU memory constraints (560M parameters)
also forced micro-batch size 2 with gradient accumulation, limiting effective per-step
gradient quality.

\subsection{Information Retrieval Module}
\label{subsec:ir}

Three retrieval channels (BM25~\cite{robertson2009bm25}, TF-IDF, dense MiniLM embeddings) are fused by Reciprocal Rank Fusion~\cite{cormack2009reciprocal} ($\text{RRF}(d)=\sum_i 1/(k+r_i(d))$, $k=60$).
Codes with RRF score $\geq0.22\times\text{top score}$, capped at 12 per document, are predicted. The cap is critical: raising it to 25 collapses F1 from 0.662 to 0.140 (precision drops from 0.57 to 0.10 while recall rises modestly) because additional hits that pass the 0.22$\times$ score threshold are predominantly false positives.

\subsection{Dictionary Baseline}
\label{subsec:dict}

A curated CSV (1,680 term-to-code entries) is matched by an Aho--Corasick automaton with a 35-term blacklist. Cardiology procedure rules fire on keywords such as PCI/PPCI/TAVI/CABG (adding Y84, Z95, and related codes); co-occurrence rules enforce clinical consistency (e.g.\ I22~$\Rightarrow$~I21); I25~$\Rightarrow$~Z95 is applied as a recall heuristic since most I25 patients in this corpus have vascular implants, though it introduces false positives corrected downstream by merge\_and. Validation micro-F1 \textbf{0.7176} (precision 0.65, recall 0.80), covering 99.76\% of documents.

\subsection{Named Entity Recognition and Entity Linking}
\label{subsec:nerel}

Greek BERT is fine-tuned as a BIO sequence tagger with a single entity type
(\texttt{MED}) and a Partial CRF~\cite{tsuboi2008partial}. The coarse type maximises
recall; the Partial CRF enforces label consistency via Viterbi decoding and enables
learning from documents with only document-level annotations.
Training: 6 epochs, batch 8, LR $2\times10^{-5}$, early stopping patience 2.

NER spans are augmented by Aho--Corasick dictionary matching; containment-based
deduplication retains the longest span. Each mention $m$ is linked to ICD-10 code $c$
by:
\begin{equation}
\text{score}(m,c) = 0.6\cdot P_{\text{prior}}(c\mid m)
                  + 0.4\cdot\cos(\mathbf{e}_{\text{ctx}},\mathbf{e}_c)
\end{equation}
where $P_{\text{prior}}$ is the (mention, code) co-occurrence frequency from training
and $\mathbf{e}_{\text{ctx}}$, $\mathbf{e}_c$ are MiniLM embeddings of a
$\pm$200-char context window and the Greek ICD-10 description respectively.
Full document-level dictionary boosting is applied after mention linking.
Best validation micro-F1: \textbf{0.7223} at epoch~5.

\subsection{Metaheuristic Ensemble}
\label{subsec:ensemble}

Rather than fixing a single fusion rule~\cite{dietterich2000ensemble}, we define a
declarative strategy space in YAML and evaluate all strategies on the validation set,
selecting the best-scoring configuration without code changes.

\subsubsection{Base strategies.}
Eleven base strategies cover the main fusion paradigms. The primary strategy is the
weighted ensemble, a convex combination of all component probability vectors:
\begin{equation}
\hat{y} = \sum_{i=1}^{6} w_i\,\hat{y}^{(i)}, \quad w \in \mathcal{W} = \Bigl\{w \in \mathbb{R}^6 : \textstyle\sum_i w_i = 1,\; w_i \geq 0\Bigr\}
\end{equation}
Additional strategies address complementary failure modes: majority vote pools predictions
across all weight-search restarts (both classic and VNS); gated secondary falls back to
IR/NER for documents where the weighted ensemble is uncertain; top-K pruning caps
predictions per document; frequency-bucketed thresholding applies separate thresholds
per label-frequency band; per-label champion routes each label to its highest-validation-F1
component; per-label champion + vote adds a label if the champion fires or a minimum
number of other models agree; per-patient score routing selects the best-scoring model
per document; and label correction mines co-occurrence rules $(a \!\to\! b)$ from
validation predictions, predicting $b$ whenever $a$ is predicted above a mined threshold.

\subsubsection{Composition operators.}
Three label-set algebra operators defined in \texttt{strategy\_compositions.yaml} compose
pairs of base strategy outputs $\{\hat{L}^{(s)}\}_{s \in \mathcal{S}}$, where
$\hat{L}^{(s)} = \{l : \hat{y}_l^{(s)} \geq t_l\}$:
\begin{align}
\hat{L}^{\text{or}}  &= \bigcup_{s \in \mathcal{S}} \hat{L}^{(s)}
  & &\text{(maximises recall)} \\
\hat{L}^{\text{and}} &= \bigcap_{s \in \mathcal{S}} \hat{L}^{(s)}
  & &\text{(maximises precision)} \\
\hat{L}^{k\text{-of-}n} &= \Bigl\{l : \sum_{s \in \mathcal{S}} \mathbb{1}[l \in \hat{L}^{(s)}] \geq k\Bigr\}
  & &\text{(balances precision/recall)}
\end{align}
The best submission applies merge\_and over the weighted ensemble and the label
correction module, letting the correction step remove false positives from the
ensemble output.

\subsubsection{Weight optimisation.}
Weights are optimised directly on the non-differentiable validation micro-F1 objective:
\begin{equation}
w^{*} = \underset{w\,\in\,\mathcal{W}}{\arg\max}\; \text{F1}_{\mu}\!\left(\hat{y}(w),\,y^{(\text{val})}\right)
\end{equation}
Two optimisers run independently, each for 2 restarts of 10,000 evaluations:
(1) \textit{classic}: random search over $\mathcal{W}$ followed by pairwise hill-climbing;
(2) \textit{VNS}~\cite{mladenovic1997variable}: shaking (radius-growing perturbation,
$k_{\max}{=}5$) followed by a short hill-climb, with $k$ reset on improvement and a
random restart when all neighbourhoods are exhausted.
The optimiser achieving higher validation micro-F1 provides the final weights;
restart predictions from both are pooled for majority-vote strategies.

\subsubsection{Submitted strategies.}
The five submissions span the precision--recall trade-off: standalone Greek BERT
(80/10 split and 100\% data) establish the single-model ceiling; three ensemble
submissions apply merge\_and (precision-oriented), merge\_k2 (balanced), and a
safer-threshold merge\_k2 variant. The merge\_and composition achieves the best F1
by combining the weighted ensemble with the label correction module's precision boost.
